\title{Large Language Models Can Automatically Engineer Features for Few-Shot Tabular Learning}

\begin{abstract}
Large Language Models (LLMs) have exhibited significant proficiency in addressing complex, novel reasoning challenges, suggesting that they could be invaluable for tabular learning—a domain critical to numerous practical applications. We introduce a new in-context learning framework, \textsf{FeatLLM}, which leverages LLMs as feature generators to construct an enhanced dataset tailored for tabular prediction tasks. These synthesized features are subsequently utilized by straightforward downstream machine learning algorithms—such as linear regression—which achieve strong few-shot learning performance. Uniquely, \textsf{FeatLLM} operates solely with these simplified predictive models and the features identified during inference, negating the necessity of querying the LLM for each new instance. In addition, \textsf{FeatLLM} only requires access to LLMs via API calls, thus bypassing constraints associated with prompt length. Through extensive evaluation on diverse tabular datasets spanning multiple domains, \textsf{FeatLLM} is shown to produce robust rules that outperform competing methods such as TabLLM and STUNT by a significant margin (on average, 10\%).
\end{abstract}

\section{Introduction}
Large Language Models (LLMs) have recently demonstrated strong capabilities in generating relevant outputs for previously unseen assignments, all without the need for task-specific model adjustment~\cite{creswell2022selection,imani2023mathprompter,taylor2022galactica}. Their exposure to vast textual corpora endows LLMs with extensive background knowledge, making them effective stand-ins for specialized expertise across various domains~\cite{Genomes2021,singhal2023large,wilcox2003role}. This capacity is instrumental for automating a spectrum of tasks, from dialogue assessment~\cite{finch2023leveraging} to code correction~\cite{fan2023automated}. By integrating descriptions of the problem and a handful of examples into their prompts, LLMs are capable of discerning task context and directing their attention toward pertinent features and data points~\cite{brown2020language,zhao2023survey}, illustrating their competency in data interpretation and implying their suitability as automated feature engineers.

Efforts to harness the domain knowledge and reasoning aptitude of LLMs have extended into tabular learning tasks. As an example, incorporating background information—such as the higher prevalence of specific diseases among older adults—can enhance disease prediction accuracy. Contemporary studies achieve this by phrasing task and feature explanations in natural language, converting tabular data accordingly, and providing this input to an LLM~\cite{dinh2022lift,wang2023anypredict}. These methods often employ in-context learning~\cite{nam2023semi} or lightweight fine-tuning approaches~\cite{dinh2022lift,hegselmann2023tabllm}. The domain-relevant knowledge encoded within LLMs has proven valuable in these contexts, surpassing standard tabular learning approaches, particularly when data is sparse~\cite{hegselmann2023tabllm}.

Existing approaches to tabular learning using LLMs face several drawbacks. Performing end-to-end predictions generally demands at least one inference from the LLM for each data point, leading to significant computational overhead. High performance frequently hinges on fine-tuning the LLM, which is impractical for models lacking openly available training frameworks or tuning interfaces—a growing issue, since the latest state-of-the-art LLMs often only offer restricted access through inference APIs~\cite{anil2023palm,openai2023gpt4}. Additionally, most of these methods are ill-equipped for dealing with long input prompts~\cite{liu2023lost}; as tabular data scales in feature count, resulting in lengthy serialized inputs, the effectiveness of such methods commonly declines or becomes completely unmanageable. These issues fundamentally restrict the practical adoption of LLM-based tabular learning strategies in operational contexts.

To overcome these challenges, our method reimagines the role of LLMs—not as engines for end-to-end inference, but as tools for feature construction that exploit their encoded background knowledge more efficiently. Concretely, we present \textsf{FeatLLM}, a new in-context learning paradigm that seeks to uncover the "criteria" or rationales employed by LLMs during prediction. For example, in a disease classification problem, the LLM can derive and articulate explicit rules correlating feature patterns with diagnoses. These criteria become the basis for engineered features, effectively substituting or enhancing the original tabular variables for downstream tasks. By repositioning the LLM to focus on feature extraction, we make it possible to utilize simpler downstream predictors; after suitable features have been generated, complex inference engines are no longer a necessity, leading to faster inference times. This framework capitalizes on LLMs' in-context learning capabilities by generating features through strategic prompt construction involving examples, obviating the need for additional model training. Moreover, the adoption of ensemble techniques—such as feature bagging~\cite{breiman2001random}—enables us to address the prompt length issue, further enhancing scalability and robustness.

\textsf{FeatLLM} decomposes the prompt design for feature engineering into two main components: (1) comprehending the task at hand and deducing connections between input features and target outcomes; and (2) leveraging both the insights from the initial phase and a small number of example cases to formulate discriminative rules that enable class separation. This stepwise approach guides LLMs in disregarding irrelevant features during rule induction. The resulting rules are executed over the dataset—interpreted as program logic—resulting in binary-valued features that capture whether individual samples comply with specific rules. Subsequently, a simple predictive model is trained to estimate class membership probabilities using these derived features. Model robustness is enhanced using an ensemble methodology that repeatedly regenerates features while employing feature bagging; by finely tuning the selection of both features and samples during bagging, \textsf{FeatLLM} alleviates issues related to the potential inflation of prompt length. Free from any requirements for model retraining, \textsf{FeatLLM} incurs minimal inference costs, making it suitable for real-world scenarios.

We assess the performance of \textsf{FeatLLM} on 13 tabular datasets under limited data conditions, demonstrating notable accuracy and consistency. Our findings highlight that LLMs are effective in exploiting domain knowledge and dataset-derived information to create high-value features. When compared to existing few-shot learning approaches, our framework consistently exhibits superior results across a variety of settings. Source code is provided at an anonymized GitHub repository:~\texttt{https://github.com/Sungwon-Han/FeatLLM}.

\section{Related Work}

\subsection{Few-Shot Learning with Tabular Data} The emergence of few-shot learning has allowed for the derivation of broadly applicable representations from datasets, even when labeled instances are limited~\cite{chen2018closer}. While the initial progress of few-shot learning was centered on image-based tasks, this paradigm has since demonstrated strong adaptability to a variety of data forms, encompassing text, audio, and, more recently, tabular datasets~\cite{majumder2022few,nam2023stunt,schick2021exploiting}. However, relying on scant labeled data may result in the acquisition of spurious associations~\cite{bartlett2021deep}. To counteract this, contemporary approaches incorporate supplementary data resources or inject domain-specific inductive biases, thereby guiding the model towards more meaningful learning. Strategies include leveraging copious, easily accessible unlabeled data with carefully crafted augmentation procedures for semi-supervised learning~\cite{ucar2021subtab,yoon2020vime}, or obtaining optimal initial parameters through unsupervised pre-training agnostic to specific tasks~\cite{somepalli2021saint}. Unsupervised training frameworks often involve introducing noise—via data masking or corruption—followed by reconstruction objectives~\cite{arik2021tabnet,majmundar2022met}, or deploying data augmentation to facilitate contrastive learning methods~\cite{bahri2022scarf,chen2020simple}. Nonetheless, because of the inherent heterogeneity of tabular data, it remains difficult to design augmentation schemes that are broadly effective, limiting the success of these methods in few-shot contexts~\cite{nam2023stunt}.

A more recent trend involves constructing models by utilizing a wide assortment of tabular datasets that extend beyond the immediate target domain, subsequently employing transfer learning to adapt to novel tasks~\cite{zhang2023generative,levin2022transfer,zhu2023xtab}. For example, TransTab~\cite{wang2022transtab} leverages transformers trained not just on raw table entries but also on semantic cues derived from column names. By extracting information from a diversity of tables and column structures, these models enable few-shot or even zero-shot inference capabilities in unfamiliar tabular learning scenarios. On the other hand, TabPFN~\cite{hollmann2022tabpfn} uses mixtures of statistical distributions to synthesize training data for pre-training Bayesian neural networks. After this pre-training step, inference involves simultaneously processing both support (training) and query (test) samples for the target problem, requiring no further learning. Accumulating knowledge from different datasets in this manner enables such models to surpass classic tree-based techniques, particularly within few-shot learning regimes.

\subsection{Language-Interfaced Tabular Learning} Employing large language models (LLMs) offers a promising direction to integrate prior knowledge into tabular learning~\cite{anil2023palm,openai2023gpt4}. Owing to their extensive pre-training on large and diverse text corpora, LLMs are capable of transferring their broad knowledge base to previously unseen tasks across multiple domains~\cite{brown2020language}. Recent investigations have begun to exploit this latent knowledge in LLMs for handling tabular data. The standard methodology involves transforming tabular structures into sequential textual representations, which are subsequently incorporated into LLM input prompts~\cite{dinh2022lift,hegselmann2023tabllm,wang2023anypredict}. By embedding not only the raw tabular entries but also feature and task explanations within these prompts, models can discern relevant feature-task interactions, thereby supporting prediction or inference. Developments in this sphere include fine-tuning LLMs with methods tailored for efficient parameter updates, such as LoRA~\cite{hu2021lora} or IA3~\cite{liu2022few}, as well as adopting in-context learning strategies that provide a handful of labeled examples within the prompt, enabling the model to reason without additional fine-tuning~\cite{dinh2022lift,hegselmann2023tabllm,nam2023semi,slack2023tablet}.

Although the previously mentioned approaches have demonstrated success in enhancing few-shot tabular learning, the necessity for all predictions to go through an LLM during inference introduces significant difficulties, particularly in industrial applications. The present work seeks to move beyond this traditional black-box framework, where every input sample must be processed end-to-end by an LLM for prediction. Rather, it proposes directing the LLM to deduce explicit rules or conditions that characterize each output class. \textsf{FeatLLM} transforms these identified rules into executable program code, facilitating the generation of novel features. As a result, subsequent predictions can be performed independently of the LLM, benefiting from in-context learning for rule derivation from both prior knowledge and limited labeled examples.

\section{Methods}
\textsf{FeatLLM} is constructed to systematically obtain "rules"—that is, the class-specific conditions—from the provided few-shot samples, as illustrated in Figure~\ref{fig:main_model}. An LLM is tasked with extracting these rules, which are then converted into new binary indicators for each sample, marking whether a given rule holds. These constructed binary features are subsequently combined through a weighted dot product, where the weights are constrained to be non-negative and are learned during training, to produce a probabilistic estimate for each class. The model thus learns the relative significance of each binary feature with respect to class prediction (see Section~\ref{Sec:training} for further explanation). To enhance stability and generalization, this procedure incorporates repeated bagging, and results from individual models are merged through an ensemble approach. The subsequent sections describe the formal problem setup as well as specifics for each stage of the methodology.

\vspace*{-0.12in}
\paragraph{Problem formulation.} We examine a tabular dataset comprising $N$ labeled instances, denoted as $\mathcal{D} = \{(\mathbf{x}^i, \mathbf{y}^i)\}_{i=1}^{N}$. Each sample in $\mathcal{D}$ is represented by $d$ features, where every $\mathbf{x}^i$ forms a $d$-dimensional feature vector. Associated with these features are human-readable feature names $F = \{f_j\}_{j=1}^{d}$ and their respective definitions, available separately. Within the scope of classification tasks, each label $\mathbf{y}^i$ belongs to one of several predetermined categories. For $k$-shot learning scenarios, the model is trained on only $k$ labeled instances selected at random from the dataset, where $k < N$.

\subsection{Prompt Design for Extracting Rules}
\label{Sec:prompt_design}
To facilitate the LLM in deriving rules through a precise reasoning process, we structure the prompts to emulate the analytical approach of a domain expert dealing with tabular data. As detailed in Fig.~\ref{fig:prompt_for_rule} and Appendix~\ref{sec:example_rule}, the constructed prompt consists of three key segments:

\vspace*{-0.12in}
\paragraph{Basic information description.} The initial component offers all necessary context required for addressing the task (see the highlighted orange section in Fig.~\ref{fig:prompt_for_rule}). This includes an explanation of the task at hand, specifications of the labels, comprehensive feature descriptions, and illustrative examples. The task overview is phrased as a direct question (e.g., ``Does this patient's myocardial infarction complications data indicate chronic heart failure? Yes or no?''), with additional samples provided in Appendix~\ref{sec:task_desc}. Feature descriptions state the type of value for each feature and, where applicable, supply the feature’s definition; categorical features may also list the range of possible values. For exemplification, several training examples with their labels are rendered as textual demonstrations. This serialization adheres to the structure adopted in earlier works~\cite{dinh2022lift,hegselmann2023tabllm}:

\begin{align}
&\text{Serialize}(\mathbf{x}^i,\mathbf{y}^i, F) = \nonumber \\ 
&``\ f_1\ \text{is}\ \mathbf{x}^i_1.\ ... \ f_d\  \text{is}\  \mathbf{x}^i_d.\ \ \text{Answer:}\  \mathbf{y}^i\ ”
\end{align}

\vspace*{-0.12in}
\paragraph{Reasoning instruction.} Our objective is to bolster the reasoning abilities of the LLM by introducing targeted reasoning instructions (as illustrated by the blue annotations in Fig.~\ref{fig:prompt_for_rule}). The instruction starts with an opening sentence akin to the chain-of-thought methodology~\cite{wei2022chain}. Subsequently, we separate the model's inference process into two distinct phases. The initial phase directs the LLM to hypothesize about possible causal relationships or tendencies between features and the main task, explicitly without consulting concrete demonstrations. Instead, the LLM is expected to leverage general background knowledge or commonsense understanding, allowing it to apply prior domain knowledge in addressing the problem. During the second phase, the LLM synthesizes insights from the given examples together with the information gathered in the first phase to derive a set of rules tailored for each class. This bifurcated strategy helps steer the LLM away from picking up on superficial or irrelevant correlations within unnecessary columns, instead enabling concentration on features of greater importance. To prevent scenarios of underfitting (too few rules) or unmanageable prompt length (too many rules), we fix the number of rules extracted per class at 10\footnote{Analysis on the number of rules can be found in Section~\ref{sec:analysis}.}. Furthermore, when drafting these rules, the LLM is instructed to refer to the feature type information from the basic information section to ensure correct type matching in each rule formulation.

\vspace*{-0.12in}
\paragraph{Response instruction.} To streamline both the parsing and subsequent application of inferred rules, we employ explicit response instructions to standardize the structure of the LLM’s outputs (highlighted in yellow in Fig.~\ref{fig:prompt_for_rule}). The prompting protocol spells out the required response format corresponding to each answer class as well as the expected template for individual rule statements. By providing these instructions, the LLM is empowered to select formats appropriate for specific feature types. If no further directives are supplied, the model has the flexibility to compose composite rules utilizing logical connectives such as AND and OR. A sample of the resulting output can be found in Appendix~\ref{sec:outcome-rule}.

\subsection{Inferring Class Likelihood via Rules}
\label{Sec:training}

\paragraph{Parsing rules for feature generation.} In this step, we employ the previously derived rules to produce a set of new binary features for our dataset. Each binary feature corresponds to a rule linked with a specific class and signals whether a sample meets the requirements of that rule (assigning a '1' if satisfied, '0' otherwise). These features thus encode how well the sample matches the rule-based definition for each class, which is useful for downstream classification. Because these rules are articulated in natural language by the LLM, an effective parsing mechanism is essential to systematically translate them into binary indicators. Even though we set formatting guidelines to make LLM-generated responses easier to parse, deviations and unexpected syntactic variations frequently occur in practice~\cite{kovavc2023large}. For example, the LLM might use phrases such as ``greater than" or ``less than" in place of symbolic operators like ``$>$" or ``$<$", introducing further variation and complicating rule extraction.

To overcome the complexities arising from textual noise in the rules, we turn to the LLM itself instead of implementing sophisticated parsing algorithms. Owing to its capability to discern underlying semantic meaning despite surface-level variations, and because it can reliably produce executable code (having been trained on programming data), we craft prompts containing the function name, descriptions of inputs and outputs, and the rules to be interpreted. These prompts are presented to the LLM to generate parsing code, as illustrated in Figures~\ref{fig:prompt_for_parsing} and ~\ref{sec:example_parsing}-\ref{sec:example_parse_outcome} in Appendix, where both the prompt templates and example outputs are provided. The code produced is then run using Python’s \texttt{exec()} function to enable automatic conversion of samples into their binary feature representations.

\vspace*{-0.12in}
\paragraph{Inferring class likelihood.} With the binary feature vectors constructed from rule adherence for each class, we estimate the likelihood that a sample belongs to each class. A straightforward approach involves tallying the number of satisfied rules per class (i.e., summing the respective binary feature vector entries). Nevertheless, since rules can differ in informativeness, their individual contributions should be learned from training data. To capture this, we employ a traditional machine learning (ML) technique that learns weights for each binary feature associated with a class, thus reflecting their predictive significance. For instance, we might use a linear model without an intercept. Formally, let $\mathbf{z}_k^i \in \{0, 1\}^R$ denote the vector of binary features for class $k$ corresponding to sample $\mathbf{x}^i$, where $R$ counts the rules per class and $c$ the total number of classes. We then compute the class probabilities $\mathbf{p}^i$ as:

\begin{align}
\text{logit}_k^i &= \max(\mathbf{w}_k, 0) \cdot \mathbf{z}_k^i, \nonumber \\
\mathbf{p}^i &= \text{Softmax}({[\text{logit}_{1}^i, ..., \text{logit}_{c}^i]}). \label{eq:infer_prob}
\end{align}

In this framework, $\mathbf{w}_k$ denotes the parameter to be trained in the linear model and reflects each feature's relative influence. Constraining these weights to a positive domain guarantees that the LLM-derived features only contribute positively, avoiding any decrease in class probability during model training. After computing the logits, the softmax function is applied to transform them into class probabilities. The parameters $\mathbf{w}_k$ are optimized using cross-entropy loss, learned over a set of few-shot training examples.

\vspace*{-0.12in}
\paragraph{Ensembling with bagging.} 
To construct the ensemble, the entire modeling workflow is executed multiple times, yielding a collection of predictive models. The final class prediction aggregates the output, using the set of optimized weights from $T$ runs, i.e., $\{\mathbf{w}[t]\}_{t=1}^T$. For each trial, a class probability vector $\mathbf{p}^i$ is computed, and these probabilities are averaged as outlined in Eq.~\ref{eq:infer_prob}.

Randomness is introduced into rule generation for each ensemble member by employing a sufficiently elevated temperature during LLM inference. Furthermore, the order of few-shot exemplars is shuffled, inspired by empirical findings that demonstration order can affect LLM responses~\cite{lu2022fantastically}\footnote{See Appendix~\ref{sec:diversity} for analyses of rule diversity.}. Bagging is then used to further leverage predictive diversity and fortify model robustness, especially when the training set or feature set is large; subsets of features or data instances are chosen for each individual model run. This ensemble strategy imparts two key benefits. First, errors arising from occasional incorrect rule generation in certain ensemble members can be mitigated by the presence of correct inferences from other members, thereby improving system resilience—a phenomenon linked to the LLM's self-consistency, whereby rules frequently produced across multiple trials are more likely to be valid~\cite{wang2022self}. Second, ensembling addresses constraints posed by limited LLM input context lengths. If the tabular dataset exceeds what can fit within the model's prompt, bagging reduces input size and preserves reliability. In cases where single-model rule sets cannot encapsulate all feature relationships, aggregating across diverse weak learners from multiple ensemble trials compensates for the limited scope, thus maintaining high overall predictive accuracy.

\section{Experiments}
We assess the performance of \textsf{FeatLLM} on a diverse set of tabular datasets and perform a range of analyses to elucidate the behavior of our approach. Due to space limitations, supplementary experiments—such as exploring alternative LLMs, qualitative evaluations, and additional studies—are provided in the Appendix.

\subsection{Performance Evaluation}
\paragraph{Datasets.}
Our evaluation employs 13 tabular datasets spanning both binary and multi-class classification problems: (1) Adult~\cite{asuncion2007uci}, predicting whether an individual's annual income surpasses \$50,000; (2) Bank~\cite{moro2014data}, predicting a client's likelihood to subscribe to a term deposit; (3) Blood~\cite{yeh2009knowledge}, predicting donor return behavior for further donations; (4) Car~\cite{kadra2021well}, classifying car quality; (5) Communities~\cite{misc_communities_and_crime_183}, estimating regional crime rates; (6) Credit-g~\cite{kadra2021well}, identifying whether an applicant presents a favorable or unfavorable credit risk; (7) Diabetes\footnote{\texttt{kaggle.com/datasets/uciml/pima-indians-diabetes-database}}, classifying diabetes presence; (8) Heart\footnote{\texttt{kaggle.com/datasets/fedesoriano/heart-failure-prediction}}, detecting coronary artery disease occurrence; and (9) Myocardial~\cite{misc_myocardial_infarction_complications_579}, predicting chronic heart failure incidents.

Additionally, we include both novel and synthetic datasets that have rarely been analyzed in previous research and were absent from LLM pre-training sets: (10) Cultivars, for predicting soybean cultivar grain yield\footnote{\texttt{archive.ics.uci.edu/dataset/913}}; (11) NHANES, for categorizing individuals' age groups based on their record\footnote{\texttt{archive.ics.uci.edu/dataset/887}}; (12) Sequence-type, to assign sequences to one of four classes: arithmetic, geometric, Fibonacci, or Collatz; and (13) Solution-mix, determining if a mixture of four solutions has a concentration percentage over 0.5. The Cultivars and NHANES datasets were added to the UCI Machine Learning Repository after September 2023, making it highly unlikely they appeared in the training data for the LLM API (gpt-3.5-turbo-0613) we utilized. Sequence-type and Solution-mix are synthetic, ensuring they could not be memorized by the language model prior to evaluation.

The selected datasets encompass a broad spectrum of sizes and complexity, as shown in Table~\ref{Tab:basic-desc}. Each includes well-defined names and attribute descriptions. Particularly, datasets such as `Communities' and `Myocardial' possess more than 100 features, presenting specific challenges due to the limited context window of LLMs.

\vspace*{-0.12in}
\paragraph{Baselines.}
Our evaluation involves comparison against nine benchmark methods. The initial three are traditional tabular modeling approaches: (1) Logistic regression (LogReg), (2) XGBoost~\cite{chen2016xgboost}, and (3) RandomForest~\cite{ho1995random}. The next two, (4) SCARF~\cite{bahri2022scarf} and (5) STUNT~\cite{nam2023stunt}, are designed to exploit unlabeled data, though, in practical scenarios, assembling extensive unlabeled tabular collections may be prohibitive. Another recent alternative is (6) TabPFN~\cite{hollmann2022tabpfn}, which relies on large-scale synthetic data generation to pretrain its model. The final group consists of LLM-driven baselines: (7) In-context learning (In-context)~\cite{wei2022emergent}, which embeds a small set of training examples directly into the LLM prompt; (8) TABLET~\cite{slack2023tablet}, which enhances prompt-based inference using auxiliary classifier-derived rule sets and prototypes; and (9) TabLLM~\cite{hegselmann2023tabllm}, which adopts efficient parameter tuning strategies, such as IA3~\cite{liu2022few}, leveraging a small number of training examples for LLM adaptation.

\vspace*{-0.12in}
\paragraph{Implementation details.}
\textsf{FeatLLM} adopts GPT-3.5 as its primary LLM, although its architecture allows for straightforward adaptation to alternative LLM choices (see Appendix~\ref{sec:backbones} for further comparative analysis). During inference, we configure the temperature parameter to 0.5 and maintain the API's default setting of 1 for top-p. The ensemble size and rule extraction count are assigned values of 20 and 10, respectively. For additional analysis regarding the sensitivity of these hyper-parameters, consult Figure~\ref{fig:hyperparameter2} and Appendix~\ref{sec:hyper-parameter}.

Model optimization employs the Adam algorithm with a learning rate of 0.01 for the linear component, with 200 epochs for training. Epochs are chosen via $k$-fold cross-validation to maximize validation performance. For the baseline reproductions, we use GPT-3.5 for in-context learning scenarios, while following the original settings for TabLLM~\cite{sanh2022multitask}, which specifies T0 as its language model. TabLLM only permits language models with openly available checkpoints, thus excluding GPT-3.5. For traditional algorithms such as LogReg, XGBoost, and RandomForest, we optimized hyper-parameters utilizing grid search in conjunction with $k$-fold cross-validation, selecting $k$ to be either 2 or 4—ensuring that every class is present within the training split at least once. Expanded descriptions of our implementation process can be found in Appendix~\ref{sec:baselines}.

\vspace*{-0.12in}
\paragraph{Results.}
Tables~\ref{tab:main1} and \ref{tab:main2} display comparative results spanning multiple datasets. Each experiment is conducted three times, varying the random split of the training and test data, and the average as well as standard deviation of the AUC metrics are reported. Across benchmarks, \textsf{FeatLLM} emerges as either the top or second-best model, outperforming the majority of baselines. Remarkably, with as few as 4 shots, our framework achieves performance akin to standard tabular methods that typically require 32 shots (refer to Fig.~\ref{fig:compares}). Although STUNT and TabPFN achieve marked improvements on select datasets, their overall performance gains over classical machine learning techniques are limited. Furthermore, LLM-based models such as those relying on in-context learning, TABLET, or TabLLM, fail to perform inference on high-dimensional tabular data. In contrast, by leveraging bagging and ensembling, our approach maintains robust performance even as the feature space grows, demonstrating scalability and effectiveness. It should be highlighted that this bagging and ensemble methodology could theoretically be integrated into current LLM-powered methods, though doing so would double per-sample inference requirements and greatly inflate computational demand, ultimately rendering such a strategy impractical.

\subsection{Analysis \& Discussion}
\label{sec:analysis}
\paragraph{Ablation study.} To investigate the influence of core framework components, we undertake a series of ablation studies focusing on: (1) \textsf{-Tuning}: removing the linear model's weight tuning and instead aggregating binary features to produce the logit; (2) \textsf{-Ensemble}: disabling the ensembling step; (3) \textsf{-Description}: excluding the incorporation of feature descriptions; and (4) \textsf{-Reasoning}: omitting the Step-1 portion of reasoning instruction during rule generation.

A summary of the average impact of these ablations on AUC across datasets is presented in Table~\ref{tab:ablation}. In every instance, excluding a component leads to decreased performance. Notably, as the number of shots rises, weight tuning confers greater benefits, highlighting that ample data enables more accurate rule importance estimation and thus augments the contribution of tuning. Conversely, when data is scarce, feature descriptions and reasoning instructions are particularly impactful, implying that LLM-driven prior knowledge becomes essential to boost results under limited supervision. Finally, the ensemble technique proposed here consistently delivers performance enhancements across all evaluated scenarios.

\vspace*{-0.12in}
\paragraph{Training \& inference time comparison.} We evaluate and contrast the training and inference durations among several techniques. All experiments utilize the Adult dataset, specifically employing a training set with 16 examples. Unless otherwise indicated, a single A100 GPU is employed, with TabLLM being the exception, leveraging four A100 GPUs to facilitate model parallelism. Table~\ref{tab:cost} details the measured runtimes. Importantly, \textsf{FeatLLM} achieves a notably low inference time that rivals that of traditional approaches. In stark contrast, other LLM-based strategies—such as In-context learning, TABLET, and TabLLM—display markedly longer inference times. With respect to training, \textsf{FeatLLM} demonstrates comparable efficiency to its LLM-based counterparts, necessitating only 30 API queries, a minimal requirement in terms of computational and financial resources, which is also amenable to parallel processing.

\vspace*{-0.12in}
\paragraph{Quality of features generated by \textsf{FeatLLM}.}
A straightforward experiment is performed to assess the informativeness of the rule-derived features output by \textsf{FeatLLM} for practical tasks. This evaluation entails calculating the average AUROC between the newly constructed features and the corresponding target labels, subsequently determining the fraction of features whose AUROC exceeds 0.5 for each examined dataset. Table~\ref{tab:quality} presents these findings. The results demonstrate that most generated rules are pertinent to the intended target attribute, thus delivering valuable information for the task (see the ``Before tuning" column). Furthermore, while tuning the linear model on the data, those rules deemed less relevant—identified by learned weights falling below a specified threshold—are systematically excluded (refer to the ``After tuning" column). The threshold value is set to 0.1 based on the overall distribution of weight magnitudes.

\vspace*{-0.12in}
\paragraph{Effect of adding spurious correlations.} We investigated how well different approaches can disregard irrelevant spurious correlations in low-data settings. For this purpose, we performed an experiment using the Heart dataset, to which we appended randomly selected columns from the Adult dataset that bear no relevance to the Heart dataset's classification objective. We then evaluated the impact on model performance as the number of extraneous columns increases. To ensure a balanced assessment, we fixed the number of labeled training instances at 8 shots and did not include unlabeled samples, thus omitting techniques such as SCARF and STUNT from our evaluation.

Figure~\ref{fig:spurious} presents a comparison of performance deterioration due to spurious correlations. In this evaluation, \textsf{FeatLLM} exhibited the smallest decline in accuracy when faced with irrelevant features, outpacing baseline approaches. This robustness is likely attributable to the mechanism within our framework that compels the alignment between features and task relevance, leveraging prior knowledge during rule extraction. As a result, rules formed on the basis of unrelated features are rarely created, leading to stable performance. In quantitative terms, rules derived from noisy columns made up only about 1-2\% of the overall set. Nevertheless, we also found that if highly related columns from other datasets are added at random, \textsf{FeatLLM} may still capture spurious associations and misinterpret the underlying causal structure between features and the task (see Appendix~\ref{sec:spurious-appendix} for further discussion).

\vspace*{-0.12in}
\paragraph{Hyper-parameter analysis}
\textsf{FeatLLM} involves two key hyper-parameters: the ensemble size and the number of rules per class extracted during each LLM query. We carried out a series of experiments to assess how these hyper-parameters affect model outcomes. The experimental results indicate that increasing the number of ensembles leads to better performance, although these improvements taper off after approximately 20 ensembles, suggesting diminishing returns (refer to Fig.~\ref{fig:appendix-hyperparameter1} in the Appendix). Regarding the quantity of rules extracted, differences in model performance were not statistically significant. Nevertheless, selecting 10 rules strikes an optimal balance between prompt size and predictive accuracy (refer to Fig.~\ref{fig:hyperparameter2}). We hypothesize that including more rules increases the dimensionality of the input and injects additional information, but this may exacerbate overfitting, especially in situations with limited data (the low-shot setting).

\section{Conclusion}
This work introduces \textsf{FeatLLM}, a novel approach that enhances LLM-driven few-shot learning for tabular datasets. Unlike standard usage where LLMs are employed solely to make predictions on individual samples, \textsf{FeatLLM} utilizes LLMs for rule and criteria generation as part of a feature engineering process. Leveraging an ensemble framework combined with bagging, \textsf{FeatLLM} minimizes inference time, operates using only API-level access (with no need for additional training), and circumvents constraints related to feature dimensionality. Experimental results demonstrate that \textsf{FeatLLM} attains robust predictive accuracy, providing an efficient and practical strategy for leveraging LLMs in real-world tabular learning scenarios.

Although currently optimized for low-shot learning contexts, \textsf{FeatLLM} paves the way for future exploration into automated feature generation with larger datasets, extending beyond rule-based features, and promoting interpretability. These enhancements aim to broaden its applicability to diverse machine learning tasks.

\section{Broader Impact} The \textsf{FeatLLM} framework is conceived to leverage the pretrained knowledge encapsulated within LLMs, thereby surpassing the standard performance limitations inherent to supervised learning on tabular datasets. By improving data efficiency and supplementing limited training samples with LLM-derived prior knowledge, our approach helps to counteract overfitting. \textsf{FeatLLM} holds particular promise in real-world industries (e.g., finance and healthcare) where annotation is expensive or logistically difficult, and where the rich, relevant knowledge in LLMs (trained on industry-specific text) may add significant value. Looking ahead, it is imperative to critically analyze potential risks—such as societal biases or misinformation inherent to LLM knowledge—that could affect predictions, especially as such biases may at times dominate over the signals present in limited labeled data, thus informing future directions for responsible and widespread deployment.

\end{document}